% GENERATED FILE. Edit canonical lessons, then run npm run generate:books.
% canonical-source-hash: fnv1a64:544210ff509ad95e
% canonical-lessons: TA-C328-mannar, TA-C328-rani, TA-C328-ilavaracan, TA-C328-ilavaraci, TA-C328-ariyanai

\chapter{Kings and Queens}
\label{ch:ta-kings-and-queens}
\hlchaptermodality{328}

\begin{chapteropening}
\textbf{By the end of this chapter:} \emph{I can talk about kings, queens and their children.}

Complete the last lesson of chapter 328: I can talk about kings, queens and their children.
\end{chapteropening}

\section[\texorpdfstring{\ta{மன்னர்} (maṉṉar)}{மன்னர் (maṉṉar)}]{\ta{மன்னர்} (maṉṉar) — a king}

\label{lesson:TA-C328-mannar}



\begin{quote}
Before the new one: say the Tamil for a long time, then the Tamil for ancestral roots.
\end{quote}

\subsection*{You'll want to know: \ta{மன்னர்}}
\textbf{\ta{மன்னர்}} — \emph{maṉṉar} — \textquotedblleft{}a king\textquotedblright{}.

\begin{quote}
\ta{மன்னர்} \ta{தன்} \ta{மக்களை} \ta{நன்றாகக்} \ta{கவனித்துக்கொண்டார்}.
\end{quote}

\begin{quote}
\emph{maṉṉar taṉ makkaḷai naṉṟākak kavaṉittukkoṇṭār.}
\end{quote}

\begin{quote}
The king took good care of his people.
\end{quote}

A respectful form; \ta{மன்னன்} is the plain one.

\begin{cousinweb}
\textbf{\ta{மன்னர்}} (\emph{maṉṉar}) is native Dravidian: DEDR 4774 gives \ta{மன்} and \ta{மன்னன்}, \textquotedblleft{}king, lord, chief\textquotedblright{}, and \ta{மன்னர்} is its respectful form.
\end{cousinweb}

\begin{grammarlens}[title={What you've built}]
One more word for telling what happened.
\end{grammarlens}

\subsection*{Your turn}
\begin{itemize}
\raggedright
  \item \emph{Say it:} \emph{maṉṉar}
  \item \emph{Say it:} \emph{maṉṉar}, once more
  \item \emph{Say it:} \emph{maṉṉar taṉ makkaḷai naṉṟākak kavaṉittukkoṇṭār.} — The king took good care of his people.
  \item \emph{Recall:} say the Tamil for a long time, then the Tamil for ancestral roots, then say \emph{maṉṉar} again
\end{itemize}

\begin{tcolorbox}[breakable,colback=teal!4,colframe=teal!35!black,title={Before you move on}]
What does \ta{மன்னர்} mean? (\textquotedblleft{}A king\textquotedblright{}.) Say it once more.
\end{tcolorbox}
\section[\texorpdfstring{\ta{ராணி} (rāṇi)}{ராணி (rāṇi)}]{\ta{ராணி} (rāṇi) — a queen}

\label{lesson:TA-C328-rani}



\begin{quote}
Before the new one: say the Tamil for ancestral roots, then the Tamil for a king.
\end{quote}

\subsection*{You'll want to know: \ta{ராணி}}
\textbf{\ta{ராணி}} — \emph{rāṇi} — \textquotedblleft{}a queen\textquotedblright{}.

\begin{quote}
\ta{ராணி} \ta{அழகான} \ta{தோட்டத்தில்} \ta{நடந்தாள்}.
\end{quote}

\begin{quote}
\emph{rāṇi aḻakāṉa tōṭṭattil naṭantāḷ.}
\end{quote}

\begin{quote}
The queen walked in a beautiful garden.
\end{quote}

Also the queen in chess and in cards.

\begin{cousinweb}
\textbf{\ta{ராணி}} (\emph{rāṇi}) is borrowed from Sanskrit \emph{rājñī}, \textquotedblleft{}queen\textquotedblright{} (CDIAL 10692), through a Prakrit form.
\end{cousinweb}

\begin{grammarlens}[title={What you've built}]
One more word for telling what happened.
\end{grammarlens}

\subsection*{Your turn}
\begin{itemize}
\raggedright
  \item \emph{Say it:} \emph{rāṇi}
  \item \emph{Say it:} \emph{rāṇi}, once more
  \item \emph{Say it:} \emph{rāṇi aḻakāṉa tōṭṭattil naṭantāḷ.} — The queen walked in a beautiful garden.
  \item \emph{Recall:} say the Tamil for ancestral roots, then the Tamil for a king, then say \emph{rāṇi} again
\end{itemize}

\begin{tcolorbox}[breakable,colback=teal!4,colframe=teal!35!black,title={Before you move on}]
What does \ta{ராணி} mean? (\textquotedblleft{}A queen\textquotedblright{}.) Say it once more.
\end{tcolorbox}
\section[\texorpdfstring{\ta{இளவரசன்} (iḷavaracaṉ)}{இளவரசன் (iḷavaracaṉ)}]{\ta{இளவரசன்} (iḷavaracaṉ) — a prince}

\label{lesson:TA-C328-ilavaracan}



\begin{quote}
Before the new one: say the Tamil for a king, then the Tamil for a queen.
\end{quote}

\subsection*{You'll want to know: \ta{இளவரசன்}}
\textbf{\ta{இளவரசன்}} — \emph{iḷavaracaṉ} — \textquotedblleft{}a prince\textquotedblright{}.

\begin{quote}
\ta{இளவரசன்} \ta{குதிரையில்} \ta{காட்டுக்குப்} \ta{போனான்}.
\end{quote}

\begin{quote}
\emph{iḷavaracaṉ kutiraiyil kāṭṭukkup pōṉāṉ.}
\end{quote}

\begin{quote}
The prince rode to the forest on a horse.
\end{quote}

\ta{இள}, young + \ta{அரசன்}, king.

\begin{cousinweb}
\textbf{\ta{இளவரசன்}} (\emph{iḷavaracaṉ}) is a Tamil compound built from \ta{இள}, \textquotedblleft{}young\textquotedblright{} (DEDR 513), and \ta{அரசன்}, \textquotedblleft{}king\textquotedblright{}, which DEDR 201 traces to Sanskrit \emph{rājan}.
\end{cousinweb}

\begin{grammarlens}[title={What you've built}]
One more word for telling what happened.
\end{grammarlens}

\subsection*{Your turn}
\begin{itemize}
\raggedright
  \item \emph{Say it:} \emph{iḷavaracaṉ}
  \item \emph{Say it:} \emph{iḷavaracaṉ}, once more
  \item \emph{Say it:} \emph{iḷavaracaṉ kutiraiyil kāṭṭukkup pōṉāṉ.} — The prince rode to the forest on a horse.
  \item \emph{Recall:} say the Tamil for a king, then the Tamil for a queen, then say \emph{iḷavaracaṉ} again
\end{itemize}

\begin{tcolorbox}[breakable,colback=teal!4,colframe=teal!35!black,title={Before you move on}]
What does \ta{இளவரசன்} mean? (\textquotedblleft{}A prince\textquotedblright{}.) Say it once more.
\end{tcolorbox}
\section[\texorpdfstring{\ta{இளவரசி} (iḷavaraci)}{இளவரசி (iḷavaraci)}]{\ta{இளவரசி} (iḷavaraci) — a princess}

\label{lesson:TA-C328-ilavaraci}



\begin{quote}
Before the new one: say the Tamil for a queen, then the Tamil for a prince.
\end{quote}

\subsection*{You'll want to know: \ta{இளவரசி}}
\textbf{\ta{இளவரசி}} — \emph{iḷavaraci} — \textquotedblleft{}a princess\textquotedblright{}.

\begin{quote}
\ta{இளவரசி} \ta{அரண்மனையில்} \ta{வாழ்ந்தாள்}.
\end{quote}

\begin{quote}
\emph{iḷavaraci araṇmaṉaiyil vāḻntāḷ.}
\end{quote}

\begin{quote}
The princess lived in the palace.
\end{quote}

The -\ta{இ} marks a woman, as in \ta{அரசன்} and \ta{அரசி}.

\begin{cousinweb}
\textbf{\ta{இளவரசி}} (\emph{iḷavaraci}) is a Tamil compound built from \ta{இள}, \textquotedblleft{}young\textquotedblright{} (DEDR 513), and \ta{அரசி}, \textquotedblleft{}queen\textquotedblright{}, the feminine of \ta{அரசன்}, which DEDR 201 traces to Sanskrit \emph{rājan}.
\end{cousinweb}

\begin{grammarlens}[title={What you've built}]
One more word for telling what happened.
\end{grammarlens}

\subsection*{Your turn}
\begin{itemize}
\raggedright
  \item \emph{Say it:} \emph{iḷavaraci}
  \item \emph{Say it:} \emph{iḷavaraci}, once more
  \item \emph{Say it:} \emph{iḷavaraci araṇmaṉaiyil vāḻntāḷ.} — The princess lived in the palace.
  \item \emph{Recall:} say the Tamil for a queen, then the Tamil for a prince, then say \emph{iḷavaraci} again
\end{itemize}

\begin{tcolorbox}[breakable,colback=teal!4,colframe=teal!35!black,title={Before you move on}]
What does \ta{இளவரசி} mean? (\textquotedblleft{}A princess\textquotedblright{}.) Say it once more.
\end{tcolorbox}
\section[\texorpdfstring{\ta{அரியணை} (ariyaṇai)}{அரியணை (ariyaṇai)}]{\ta{அரியணை} (ariyaṇai) — a throne}

\label{lesson:TA-C328-ariyanai}



\begin{quote}
Before the new one: say the Tamil for a prince, then the Tamil for a princess.
\end{quote}

\subsection*{You'll want to know: \ta{அரியணை}}
\textbf{\ta{அரியணை}} — \emph{ariyaṇai} — \textquotedblleft{}a throne\textquotedblright{}.

\begin{quote}
\ta{புதிய} \ta{மன்னர்} \ta{அரியணையில்} \ta{அமர்ந்தார்}.
\end{quote}

\begin{quote}
\emph{putiya maṉṉar ariyaṇaiyil amarntār.}
\end{quote}

\begin{quote}
The new king sat on the throne.
\end{quote}

\ta{அரியணை} \ta{ஏறு}, to come to the throne.

\begin{cousinweb}
\textbf{\ta{அரியணை}} (\emph{ariyaṇai}) is a Tamil compound built from \ta{அரி}, \textquotedblleft{}lion\textquotedblright{}, and \ta{அணை}, \textquotedblleft{}support, prop\textquotedblright{} (DEDR 123): a lion seat.
\end{cousinweb}

\begin{grammarlens}[title={What you've built}]
One more word for telling what happened.
\end{grammarlens}

\subsection*{Your turn}
\begin{itemize}
\raggedright
  \item \emph{Say it:} \emph{ariyaṇai}
  \item \emph{Say it:} \emph{ariyaṇai}, once more
  \item \emph{Say it:} \emph{putiya maṉṉar ariyaṇaiyil amarntār.} — The new king sat on the throne.
  \item \emph{Recall:} say the Tamil for a prince, then the Tamil for a princess, then say \emph{ariyaṇai} again
\end{itemize}

\begin{tcolorbox}[breakable,colback=teal!4,colframe=teal!35!black,title={Before you move on}]
What does \ta{அரியணை} mean? (\textquotedblleft{}A throne\textquotedblright{}.) Say it once more.
\end{tcolorbox}
